\par\noindent\textbf{(1)} 无阻力自静止下滑由能量守恒有 v=√(2gy)，飞行时间 T[y]=∫√(1+y'²)/(√(2gy)) dx。被积函数不显含 x，Beltrami恒等式 F−y'F\_\{y'\}=1/[√(2gy)√(1+y'²)]=常数，故 y(1+y'²)=2a，为摆线一阶积分。
\par\noindent\textbf{(2)} 积分 y(1+y'²)=2a，令 y=a(1−cosθ)、x=a(θ−sinθ)，可直接求 y'/x'=cot(θ/2)，于是 y(1+cot²(θ/2))=2a。沿摆线 ds/dθ=2a sin(θ/2)，而 v=2√(ga)sin(θ/2)，从 θ=0 到 θ\_f 的时间为 √(a/g)θ\_f，尽管起点速度为零、斜率发散，时间仍有限。
\par\noindent\textbf{(3)} 对任意固定光滑轨道用弧长 s，令 w=v²，则 m v dv/ds=mg dy/ds−ηv²，即 (m/2)dw/ds+ηw=mg dy/ds。初始 w(0)=0，解为 w(s)=2g e\textasciicircum{}\{−2ηs/m\}∫\_0\textasciicircum{}s e\textasciicircum{}\{2ηu/m\}(dy/du)du，需保持 w≥0 才能继续滑动。
\par\noindent\textbf{(4)} 有阻力时速度取决于沿此前弧段累积的耗散，时间泛函不再仅由局部 y,y' 表示，因此用无耗散Beltrami恒等式直接宣称摆线仍最优是无效的。令 η→0 可用积分核回收 w(s)=2gy(s)（初始y=0），从而恢复无耗散轨迹的时间泛函。
